\section{Discussion and Limitations}
\label{sec:discussion}

\subsection{Predictive Value Versus Interpretability}

\tabicl did not justify its complexity. Its advantage over XGBoost and the main-effects EBM lay
within sampling noise on every metric, and EBM matched or slightly exceeded it at the operating
point (recall 0.512 vs.\ 0.507; precision 1.000 vs.\ 0.991) while being transparent by
construction. In line with \citet{rudin2019stop}, the EBM (or logistic regression) would, on this
evidence, be the preferable choice for any further development: it offers statistically
indistinguishable performance, explanations that are the model itself rather than an approximation
of it, no sampling instability and far cheaper inference (0.007 s versus 10.8 s to score the 545
test students on a CPU).
\subsection{Interpreting the Near-Perfect Precision}

The near-perfect precision does not indicate an excellent screening tool; it arises because the top
20\% is filled with students whose records already show disengagement (Section~\ref{sec:who}). Such
students are easy to rank, yet many may already be beyond the reach of a meeting and could be found
with a transcript query; even among students who sat an evaluation, the flagged ones are those with
weak results (Section~\ref{sec:sensitivity}). For an officer, the list adds an ordering, not new
information. The more consequential group is the 105 unflagged validation dropouts, 93 of
whom passed at least one unit. A simple rule could handle the obvious cases, reserving meetings for
less visible students.

\subsection{Implications of the Explanation Audit}

The explanations are faithful and stable under small data changes. However, the exact list of top reasons often changes when the
explanation is recomputed with a different seed or comparison group, so the bars should be read as
an indication of what broadly mattered rather than as a precise ranking. Seed instability may partly
reflect the one-permutation budget, the minimum that shap allows (a larger budget was not tested),
whereas background instability reflects a substantive choice of which students count as
``typical'' and would not disappear with more computation. The local research interface
(Appendix~\ref{sec:app-interface}) displays each student's stability alongside the reasons.

\textbf{Fee status and recourse.} Tuition fees up to date was the third most important group, and
for 17 flagged students clearing a fee or debt record would remove the flag. In the model, the arrow
runs from settled fees to a lower score, but this must not be read as ``paying fees prevents
dropout'', since the causal arrow may point the other way: students
already planning to leave, or in financial difficulty, may stop paying, making unpaid fees a symptom
rather than a cause. A score change is not a causal effect~\citep{karimi2021recourse};
an unpaid fee is better treated as a prompt for a conversation.

\subsection{Interpreting the Human Study}

The human study tested whether people could predict the model's flag, not whether explanations
improve support decisions. For Group~B the task partly became a chart-reading exercise, since adding
up the bars largely reconstructs the flag (Section~\ref{sec:audit}), and it was also easy without
charts, as flagged profiles were visibly weak. The lack of an accuracy gain is therefore not proof that
explanations lack value. Nevertheless, case C01 shows that bars displayed without the cohort cut-off can make an
unflagged student appear flagged, and the higher confidence of the chart group, without any gain in
accuracy, signals a risk of over-reliance.
\subsection{Protected Attributes, Fairness and Deployment Boundaries}

Whether gender, age or nationality should influence who receives a meeting is contested. Retaining
them reflects real differences in base rates, and a model blind to them can still be unfair through
proxies such as application route and course. On validation, however, removing them cost almost
nothing (ROC-AUC 0.952 to 0.950), altered only 5 of 108 flags and narrowed the gender recall gap
from 0.19 to 0.10; moreover, allocating
outreach on the basis of who a student is, rather than how the student is performing, would be
difficult to justify. The position adopted here is therefore to \emph{audit} with protected
attributes but not to \emph{prioritise} with them. Subgroup results are descriptive, not a fairness
certificate.

\subsection{Threats to Validity}

The evidence comes from one Portuguese institution without external or temporal validation, tests
one foundation model (TabPFN was not compared), assumes rather than measures the 20\% capacity,
excludes enrolled students and rests on a single test set of 545 students whose intervals are wide relative to
model differences. The full explanation audit followed the final test, using a refit that matches
the frozen model, and its results depend on the SHAP budget and the 25-row background. The human
study used volunteers, ten selected cases and an unsupervised online format; its comparisons are
exploratory, and its coding is a single pass without measured agreement. None of the results identify causes of dropout or effects of
interventions.

